\documentclass[10pt,a4paper]{amsart}
\usepackage{amsmath,amssymb,amsthm,graphicx}
\usepackage{verbatim}
\newenvironment{reference}{\comment}{\endcomment}
%\newenvironment{reference}{}{}
\newenvironment{slogan}{\comment}{\endcomment}
\newenvironment{history}{\comment}{\endcomment}

% For commutative diagrams we use Xy-pic
\usepackage[all]{xy}

% We use 2cell for 2-commutative diagrams.
\xyoption{2cell}
% Theorem environments.
%
\theoremstyle{plain}
\newtheorem{theorem}[subsection]{Theorem}
\newtheorem{proposition}[subsection]{Proposition}
\newtheorem{lemma}[subsection]{Lemma}

\theoremstyle{definition}
\newtheorem{definition}[subsection]{Definition}
\newtheorem{example}[subsection]{Example}
\newtheorem{exercise}[subsection]{Exercise}
\newtheorem{situation}[subsection]{Situation}

\theoremstyle{remark}
\newtheorem{remark}[subsection]{Remark}
\newtheorem{remarks}[subsection]{Remarks}

\numberwithin{equation}{subsection}

% Macros
%
\def\lim{\mathop{\mathrm{lim}}\nolimits}
\def\colim{\mathop{\mathrm{colim}}\nolimits}
\def\Spec{\mathop{\mathrm{Spec}}}
\def\Hom{\mathop{\mathrm{Hom}}\nolimits}
\def\Ext{\mathop{\mathrm{Ext}}\nolimits}
\def\SheafHom{\mathop{\mathcal{H}\!\mathit{om}}\nolimits}
\def\SheafExt{\mathop{\mathcal{E}\!\mathit{xt}}\nolimits}
\def\Sch{\mathit{Sch}}
\def\Mor{\mathop{\mathrm{Mor}}\nolimits}
\def\Ob{\mathop{\mathrm{Ob}}\nolimits}
\def\Sh{\mathop{\mathit{Sh}}\nolimits}
\def\NL{\mathop{N\!L}\nolimits}
\def\CH{\mathop{\mathrm{CH}}\nolimits}
\def\proetale{{pro\text{-}\acute{e}tale}}
\def\etale{{\acute{e}tale}}
\def\QCoh{\mathit{QCoh}}
\def\Ker{\mathop{\mathrm{Ker}}}
\def\Im{\mathop{\mathrm{Im}}}
\def\Coker{\mathop{\mathrm{Coker}}}
\def\Coim{\mathop{\mathrm{Coim}}}

% Boxtimes
%
\DeclareMathSymbol{\boxtimes}{\mathbin}{AMSa}{"02}

%
% Macros for moduli stacks/spaces
%
\def\QCohstack{\mathcal{QC}\!\mathit{oh}}
\def\Cohstack{\mathcal{C}\!\mathit{oh}}
\def\Spacesstack{\mathcal{S}\!\mathit{paces}}
\def\Quotfunctor{\mathrm{Quot}}
\def\Hilbfunctor{\mathrm{Hilb}}
\def\Curvesstack{\mathcal{C}\!\mathit{urves}}
\def\Polarizedstack{\mathcal{P}\!\mathit{olarized}}
\def\Complexesstack{\mathcal{C}\!\mathit{omplexes}}
% \Pic is the operator that assigns to X its picard group, usage \Pic(X)
% \Picardstack_{X/B} denotes the Picard stack of X over B
% \Picardfunctor_{X/B} denotes the Picard functor of X over B
\def\Pic{\mathop{\mathrm{Pic}}\nolimits}
\def\Picardstack{\mathcal{P}\!\mathit{ic}}
\def\Picardfunctor{\mathrm{Pic}}
\def\Deformationcategory{\mathcal{D}\!\mathit{ef}}
\usepackage{cleveref}
\crefname{originalreference}{original reference}{original references}
\Crefname{originalreference}{Original reference}{Original references}
\makeatletter
\newlabel{spaces-morphisms-definition-flat}{{03ML}{source}{}{sourcepointer-1}{}}
\newlabel{spaces-morphisms-definition-flat@cref}{{[originalreference][0][]03ML}{[1][0][]source}{}{}{}}
\makeatother
\crefname{originalreference}{original reference}{original references}
\Crefname{originalreference}{Original reference}{Original references}
\makeatletter
\newlabel{spaces-morphisms-definition-locally-finite-type}{{03XF}{source}{}{sourcepointer-2}{}}
\newlabel{spaces-morphisms-definition-locally-finite-type@cref}{{[originalreference][0][]03XF}{[1][0][]source}{}{}{}}
\makeatother
\crefname{originalreference}{original reference}{original references}
\Crefname{originalreference}{Original reference}{Original references}
\makeatletter
\newlabel{spaces-morphisms-definition-locally-finite-presentation}{{03XP}{source}{}{sourcepointer-3}{}}
\newlabel{spaces-morphisms-definition-locally-finite-presentation@cref}{{[originalreference][0][]03XP}{[1][0][]source}{}{}{}}
\makeatother
\crefname{originalreference}{original reference}{original references}
\Crefname{originalreference}{Original reference}{Original references}
\makeatletter
\newlabel{descent-lemma-curiosity}{{06NB}{source}{}{sourcepointer-4}{}}
\newlabel{descent-lemma-curiosity@cref}{{[originalreference][0][]06NB}{[1][0][]source}{}{}{}}
\makeatother
\title{A property of a composition survives along a surjection of sheaves}
\author{The Stacks Project Authors}
\date{stacks-project:spaces-descent.tex:local-snapshot-20261012; GNU FDL 1.2 or later}
\begin{document}
\maketitle
\noindent\textit{Editorial note.} All delivered mathematical text is the author's own, copied verbatim from the pinned Stacks Project chapter file spaces-descent.tex; no statement or proof step was rewritten, reordered or removed.\par\medskip
\noindent\textit{Editorial note.} The counted span cites three Morphisms of Spaces definitions and one Descent lemma, recorded in assembly.external\_references with canonical tags as navigation pointers.\par\medskip
\noindent\textit{Allowed standard prerequisites (editorial).}\begin{itemize}
\item Definitions of flat, locally finite type, locally finite presentation for algebraic-space morphisms. Stacks Project, Morphisms of Spaces, Definitions spaces-morphisms-definition-flat, spaces-morphisms-definition-locally-finite-type, spaces-morphisms-definition-locally-finite-presentation. Reduce the property to the chart morphisms \$U \textbackslash{}to W\$ and \$V \textbackslash{}to W\$.
\item Scheme case. Stacks Project, Descent, Lemma descent-lemma-curiosity. Concludes for \$V \textbackslash{}to W\$.
\end{itemize}\medskip
\ifdefined\hypertarget\hypertarget{sourcepointer-1}{}\fi
\noindent\textit{Source reference.} Original reference 03ML: sources/stacks-chow/tags.txt, source lines 4753--4753. Cross-chapter standard prerequisite cited by the author inside the counted span; it is used as a pointer and its own proof is not part of this item.\par\medskip
\ifdefined\hypertarget\hypertarget{sourcepointer-2}{}\fi
\noindent\textit{Source reference.} Original reference 03XF: sources/stacks-chow/tags.txt, source lines 5146--5146. Cross-chapter standard prerequisite cited by the author inside the counted span; it is used as a pointer and its own proof is not part of this item.\par\medskip
\ifdefined\hypertarget\hypertarget{sourcepointer-3}{}\fi
\noindent\textit{Source reference.} Original reference 03XP: sources/stacks-chow/tags.txt, source lines 5156--5156. Cross-chapter standard prerequisite cited by the author inside the counted span; it is used as a pointer and its own proof is not part of this item.\par\medskip
\ifdefined\hypertarget\hypertarget{sourcepointer-4}{}\fi
\noindent\textit{Source reference.} Original reference 06NB: sources/stacks-chow/tags.txt, source lines 8530--8530. Cross-chapter standard prerequisite cited by the author inside the counted span; it is used as a pointer and its own proof is not part of this item.\par\medskip
\begin{lemma}
\label{lemma-curiosity}
Let $S$ be a scheme. Let $X \to Y \to Z$ be morphism of algebraic spaces.
Let $P$ be one of the following properties of morphisms of algebraic spaces
over $S$:
flat, locally finite type, locally finite presentation.
Assume that $X \to Z$ has $P$ and that
$X \to Y$ is a surjection of sheaves on $(\Sch/S)_{fppf}$.
Then $Y \to Z$ is $P$.
\end{lemma}

\begin{proof}
Choose a scheme $W$ and a surjective \'etale morphism $W \to Z$.
Choose a scheme $V$ and a surjective \'etale morphism $V \to W \times_Z Y$.
Choose a scheme $U$ and a surjective \'etale morphism $U \to V \times_Y X$.
By assumption we can find an fppf covering $\{V_i \to V\}$ and
lifts $V_i \to X$ of the morphism $V_i \to Y$. Since $U \to X$ is surjective
\'etale we see that over the members of the fppf covering
$\{V_i \times_X U \to V\}$ we have lifts into $U$. Hence $U \to V$ induces
a surjection of sheaves on $(\Sch/S)_{fppf}$.
By our definition of what it means to have property $P$ for a
morphism of algebraic spaces (see
Morphisms of Spaces,
Definition \ref{spaces-morphisms-definition-flat},
Definition \ref{spaces-morphisms-definition-locally-finite-type}, and
Definition \ref{spaces-morphisms-definition-locally-finite-presentation})
we see that $U \to W$ has $P$ and we have to show $V \to W$ has $P$.
Thus we reduce the question to the case of morphisms of schemes
which is treated in
Descent, Lemma \ref{descent-lemma-curiosity}.
\end{proof}

\end{document}
